\SourceLeaf{32}\noindent\hspace*{\parindent}The formulas given under [1] as »Theorem I« can readily be seen from this.
It is very probable, though uncertain, that
\textsc{Gauss} used the complex factors $t_1$, $t_2$
to establish them.

What \textsc{Gauss} intended by forming the sequence of numbers $t'$, $t''$, $t'''$,
.\,.\,. is unspecified.

In the developments that follow, $t$, $\varphi t$,
$Ft$ are replaced by the notations $\Delta$, $D$, $\nu\nu A+\nu B+C$. The aim is
to pursue further the relationship between the two numbers $a+b\nu+c\nu\nu$ and
$\nu\nu A+\nu B+C$. According to handwritten notes by
\textsc{Schering}, these formulas can be read off directly from the equations of
»Theorem I«. For\,example, if in

\[
F(Ft) = (\varphi t)\,t
\]

\noindent we substitute the current meanings of $t$, $\varphi t$, and $Ft$, we obtain

\[
F(Ft) = (BB-AC)\nu\nu + (nAA-BC)\nu + (CC-nAB)
\]
\[
(\varphi t)\,t = D(c\nu\nu+b\nu+a),
\]

\noindent and from the identity of the two expressions we obtain

\[
CC-nAB = Da, \qquad nAA-BC = Db, \qquad BB-AC = Dc.
\]

In \textsc{Gauss}'s handwriting, Note [2] also gives the approximate correction

\begin{center}\textit{Correction von $\nu$ à peu près} $= -\dfrac{D^2}{9BCC},$\end{center}

\noindent and the unfinished statement on a ternary form under substitution

\begin{center}\textit{Forma ternaria per subst.}\end{center}
\[
\begin{aligned}
x &= ax'+by'+cz', \\
y &= ncx'+ay'+bz', \\
z &= nbx'+ncy'+az',
\end{aligned}
\]

\Needspace{9\baselineskip}\noindent which is to be interpreted as saying that the ternary form

\[
nnx^3+ny^3+z^3-3nxyz
\]

\noindent is transformed by that substitution into

\[
D(nnx'^3+ny'^3+z'^3-3nx'y'z').
\]

The proposition [3] on cubic residues follows directly from the congruences
established by the preceding formulas.

\[
nA^3-B^3 \equiv 0, \qquad nB^3-C^3 \equiv 0, \qquad C^3-nnA^3 \equiv 0 (\mathrm{mod.}\ D).
\]

Finally, consider the two numerical tables at the end of the fragments. Apart from
the number for $n=8$, namely $-2+\sqrt[3]{64}$, and the two
entries for $n=7$ and $n=9$, which appear to be mere slips of the pen, they
give only numbers $a+b\nu+c\nu\nu$ of norm $\pm 1$,
that\,is, units. In the first table, the one for $n=2$,
$1+\sqrt[3]{2}+\sqrt[3]{4}$ is the fundamental unit of the field, whose first
nine powers are calculated along with their reciprocal units. A written
communication%
